\documentclass[11pt,a4paper]{ctexart}

\usepackage[margin=2.3cm]{geometry}
\usepackage{array}
\usepackage{booktabs}
\usepackage{enumitem}
\usepackage{hyperref}
\usepackage{longtable}
\usepackage{tabularx}
\usepackage{xcolor}
\usepackage{listings}

\hypersetup{
  colorlinks=true,
  linkcolor=blue,
  urlcolor=blue,
  pdftitle={WashMate Campus 编程大作业说明文档},
  pdfauthor={WashMate Campus Team}
}

\lstdefinestyle{command}{
  basicstyle=\ttfamily\small,
  breaklines=true,
  columns=fullflexible,
  frame=single,
  rulecolor=\color{black!30}
}

\setlist[itemize]{leftmargin=1.8em}
\setlist[enumerate]{leftmargin=2.0em}
\renewcommand{\arraystretch}{1.25}
\setcounter{tocdepth}{2}
\setcounter{secnumdepth}{2}
\sloppy

\title{WashMate Campus\\编程大作业说明文档}
\author{武子皓、王成思宇、林恺鋆、王博翔、周逸航}
\date{2026-05-31}

\begin{document}
\maketitle

\begin{center}
\begin{tabular}{rl}
课程名称： & 大模型与生成式人工智能 \\
作业名称： & 编程大作业 \\
仓库地址： & \url{https://github.com/vegetable123q/Wash-agent} \\
APK 路径： & \path{frontend/android/app/build/outputs/apk/debug/app-debug.apk} \\
\end{tabular}
\end{center}

\clearpage
\pdfbookmark[1]{目录}{toc}
\tableofcontents
\clearpage

\section{项目概述}

\subsection{文档用途}

本文档介绍 WashMate Campus 的产品目标、主要功能、前端实现、后端参考实现、数据结构、外部服务、测试验证和发布流程。它可以用于了解产品全貌，也可以作为后续维护和功能交接的说明材料。

WashMate Campus 不是“网页前端 + 长期运行的 Python HTTP 后端”的传统结构。当前稳定的用户侧产品是 \path{frontend/} 下的 React + TypeScript + Vite + Capacitor 移动端应用。核心业务服务以内置 TypeScript 模块的形式运行在 APK 内部，主要位于 \path{frontend/src/api/}。\path{backend/} 下的 Python 模块是参考实现和测试基准，用于明确模块边界、业务规则、数据结构和集成测试行为。

\subsection{问题背景与目标用户}

校园公共洗衣场景中，学生经常需要同时判断多个问题：今天是否需要洗衣，哪些衣物可以一起洗，哪些衣物不能烘干，宿舍楼下还有没有空闲机器，排队要等多久，本次洗衣大概要花多少钱，洗完后该如何晾晒或取衣。以上判断涉及衣物材质、颜色、洗护标签、穿着次数、个人偏好、机器实时状态、天气与晾晒条件等多类信息。

WashMate Campus 的目标是把这些信息放进同一个手机流程。用户维护自己的衣柜和宿舍楼配置后，应用能够给出可执行的洗衣建议：什么时候洗、洗哪些、怎么分桶、用什么机器程序、是否需要洗衣袋、是否适合烘干、预计费用和时间是多少，以及有哪些风险提醒。

\subsection{已实现的产品范围}

当前移动端已经覆盖以下页面和流程：

\begin{itemize}
  \item 今日建议（Today）：展示天气、脏衣篮状态、频率建议、方案入口和刷新状态。
  \item 穿搭 Wiki（Outfit Wiki）：基于衣柜、天气和穿搭记录生成搭配建议，并支持加入脏衣篮。
  \item 记录穿搭（Record Outfit）：记录当天穿过的衣物，并自动增加相关衣物的穿着次数。
  \item 衣柜（Wardrobe）：展示本地衣柜，支持查看详情、删除和清空。
  \item 添加衣物（Add Clothing）：支持手动录入、本地图片、ModelHub 图片识别和文字识别。
  \item 衣物详情（Clothing Detail）：展示材质、颜色、洗护记忆、风险、穿着次数和加入脏衣篮操作。
  \item 脏衣篮（Dirty Basket）：选择本次要洗的衣物，形成洗衣规划模块的输入。
  \item 洗衣房（Laundry Room）：展示宿舍楼机器状态、排队估计、价格模式和天气上下文。
  \item 机器详情（Machine Detail）：展示选中机器的用户可理解信息。
  \item 方案详情（Plan Detail）：展示分桶、程序、费用、时间和风险提醒，并支持确认执行、完成提示和撤销。
  \item 报告（Report）：展示结构化洗衣报告、现在照做的简短行动说明、本周已完成洗衣统计、风险说明和节水节电说明。
  \item 我的（Profile）：保存宿舍楼、楼层、取衣时间、烘干偏好；ModelHub 配置和演示数据入口默认收在高级设置中。
\end{itemize}

最新移动端脚本还补强了几处执行体验：Add Clothing 在识图按钮旁直接提示缺少 ModelHub 配置；Outfit Wiki 在缺少裤装或裙装时给出添加入口；Laundry Room 会展示设备编号，并高亮当前方案推荐使用的机器。完成一次洗衣后，应用会在本地保存完成记录，让报告页在脏衣篮清空后仍能显示本周统计和最近一次洗衣内容。

\section{需求与评分点对应关系}

本项目按照 \path{evaluation_system.md} 中的最终交付评测体系整理。评分点和文档、代码、测试的对应关系如下。

\begin{longtable}{p{0.18\textwidth} p{0.12\textwidth} p{0.30\textwidth} p{0.30\textwidth}}
\toprule
评分模块 & 分值 & 本项目对应内容 & 主要文件或验证方式 \\
\midrule
后端核心功能 & 25 & 衣物抽取、衣柜、校园上下文、洗衣规划和报告生成完整串联。 & \path{backend/}、\path{tests/}、完整集成测试。 \\
数据可信度 & 15 & 材质、天气、机器、价格和时长均来自明确来源；缺失时显示错误或未知状态。 & \path{config/machine_rules.json}、\path{campusMachineApi.ts}、\path{weatherService.ts}。 \\
移动端流程 & 20 & Today、Wardrobe、Add Clothing、Laundry Room、Plan Detail、Report、Profile 等主流程可用。 & \path{frontend/src/screens/}、前端集成测试、APK 冒烟测试。 \\
UI 与交互质量 & 15 & 页面有空态、加载态、错误态；按钮有明确行为或禁用原因。 & Vitest 页面测试、人工验收流程。 \\
边界与异常 & 15 & API 失败、配置缺失、空衣柜、机器不可用、天气不可用等状态显式展示。 & \path{evaluation_system.md} 边界清单、前后端异常测试。 \\
工程质量 & 10 & 测试、构建、APK、文档、CI 和模块边界完整。 & \path{README.md}、\path{docs/structure_design.md}、GitHub workflows。 \\
\bottomrule
\end{longtable}

\section{系统功能展示}

\subsection{页面截图预留}

以下位置用于放置最终页面截图。图片文件可在本地补充后替换占位框。建议截图覆盖 Today、Wardrobe、Add Clothing、Laundry Room、Plan Detail 和 Report。

\begin{center}
\begin{tabular}{ccc}
\fbox{\parbox[c][3.2cm][c]{0.28\textwidth}{\centering Today 截图预留}} &
\fbox{\parbox[c][3.2cm][c]{0.28\textwidth}{\centering Wardrobe 截图预留}} &
\fbox{\parbox[c][3.2cm][c]{0.28\textwidth}{\centering Add Clothing 截图预留}} \\
\\
\fbox{\parbox[c][3.2cm][c]{0.28\textwidth}{\centering Laundry Room 截图预留}} &
\fbox{\parbox[c][3.2cm][c]{0.28\textwidth}{\centering Plan Detail 截图预留}} &
\fbox{\parbox[c][3.2cm][c]{0.28\textwidth}{\centering Report 截图预留}} \\
\end{tabular}
\end{center}

\subsection{Demo 演示流程}

\path{docs/e_demo_script.md} 中的 3--5 分钟演示流程可以浓缩为以下路径：

\begin{enumerate}
  \item 打开移动端 App，查看 Today 页的天气、脏衣篮、洗衣建议和方案入口。
  \item 进入 Wardrobe 或 Add Clothing，查看样例衣物或添加新衣物。
  \item 在 Dirty Basket 中选择本次要洗的衣物。
  \item 进入 Laundry Room，查看宿舍楼机器状态、设备编号和排队估计。
  \item 打开 Plan Detail，查看分桶、程序、费用、时长和风险提醒。
  \item 打开 Report，查看行动步骤、本周完成记录、风险说明和节水节电说明。
\end{enumerate}

\section{技术架构与模块设计}

\subsection{技术栈}

\begin{longtable}{p{0.20\textwidth} p{0.30\textwidth} p{0.40\textwidth}}
\toprule
层级 & 技术 & 在本项目中的职责 \\
\midrule
Mobile UI & React 18, TypeScript, Vite & 手机端页面、表单、状态、导航和用户交互。 \\
APK shell & Capacitor 7, Android Gradle & Android WebView 容器、应用标识、Capacitor sync、debug/release APK 构建。 \\
APK 内业务模块 & \path{frontend/src/api/} TypeScript modules & 衣柜、脏衣篮、完成洗衣记录、校园上下文、洗衣规划、报告生成、频率建议、识图请求和本地存储。 \\
Python 参考实现 & Python 3.12, dataclasses, unittest & 参考业务实现、共享模型、规则验证和测试基准。 \\
External services & ModelHub/Gemini v1beta, CleverSchool, Haier, Open-Meteo & 衣物识别、宿舍机器状态和天气信息。 \\
Tooling & uv, npm, GitHub Actions & 依赖同步、测试、前端构建、Capacitor 同步、Android 构建和发布。 \\
\bottomrule
\end{longtable}

\subsection{核心目录}

\begin{lstlisting}[style=command]
Wash-agent/
  README.md
  AGENTS.md
  pyproject.toml
  app.py
  main.py
  backend/
    shared/models.py
    clothing_extraction/
    wardrobe/
    campus/
    laundry/
    reports/
  frontend/
    src/
      App.tsx
      api/
      components/
      screens/
      data/
    android/
    package.json
    capacitor.config.ts
  config/
    api_config.example.json
    machine_rules.json
  data/
    wardrobe_sample.json
    machines_mock.json
    pics/
  docs/
    structure_design.md
    c_delivery.md
    e_demo_script.md
    superpowers/plans/
  tests/
    test_*.py
  .github/workflows/
    preview.yml
    release-apk.yml
\end{lstlisting}

\subsection{运行时形态}

项目的实际运行时以移动端为主：

\begin{enumerate}
  \item 用户打开 Capacitor 打包的移动端应用。
  \item React 页面通过顶层 \path{frontend/src/App.tsx} 进行导航、状态编排和事件处理。
  \item 页面调用 \path{frontend/src/api/} 中的 TypeScript service，读取本地衣柜、脏衣篮、用户配置和 ModelHub 配置。
  \item 在用户配置宿舍楼后，校园上下文模块读取真实校园机器状态；天气模块读取 Open-Meteo。
  \item 洗衣规划模块根据选中衣物、用户约束、校园上下文和价格规则生成洗衣方案。
  \item report generator 把方案、完成洗衣记录和上下文转换为页面可直接渲染的报告。
  \item 用户确认执行方案后，移动端先记录本次洗衣快照，再清空脏衣篮；用户可以在成功提示中撤销。
\end{enumerate}

Python 后端不需要作为一个常驻 HTTP 服务被移动端调用。它的价值在于提供更清晰的模块化参考实现、共享业务语言和测试基准。

\subsection{架构流程图}

\begin{lstlisting}[style=command]
User
  -> React / Capacitor screens
  -> frontend/src/api/mobileSummary.ts
      -> userProfile and ModelHub config
      -> local wardrobe and dirty basket
      -> campusMachineApi + weatherService
      -> frequencyAdvisor
      -> laundryPlanner
      -> reportGenerator
      -> completed laundry records
  -> mobile screens render plan, report, machines, and wardrobe state

Python backend
  -> backend/shared/models.py defines shared data models
  -> backend/* modules define reference behavior and tests
  -> tests/ verifies extraction, wardrobe, campus, planning, report, and integration
\end{lstlisting}

\subsection{模块边界原则}

页面层只做交互和编排，不写 ModelHub 提示词、机器解析、JSON 存储细节、洗衣分桶规则和最终报告生成。跨模块数据通过明确的数据结构传递：Python 侧使用 \path{backend/shared/models.py}，移动端使用 \path{frontend/src/api/types.ts}。当已有模型能够描述数据时，不应新增平行结构或兼容旧路径。

\section{核心实现说明}

\subsection{应用入口与导航}

\path{frontend/src/App.tsx} 是移动端的应用编排中心，负责：

\begin{itemize}
  \item 当前 screen id、浏览器历史和返回逻辑。
  \item 当前用户 profile、ModelHub config 和 mobile summary。
  \item 选中衣物 id、选中机器 id、最近完成的洗衣记录和刷新状态。
  \item 顶层事件，例如保存 profile、保存或清除 ModelHub 配置、删除衣物、选择脏衣篮、完成或撤销洗衣方案、保存穿搭记录等。
\end{itemize}

应用没有引入额外路由库，而是通过 screen id 和 parent tab 映射管理页面关系。例如 \texttt{planDetail} 和 \texttt{dirtyBasket} 属于 Today，\texttt{addClothing} 和 \texttt{clothingDetail} 属于 Wardrobe，\texttt{machineDetail} 属于 Laundry Room，\texttt{recordOutfit} 属于 Outfit Wiki。

\subsection{页面职责}

\begin{longtable}{p{0.30\textwidth} p{0.60\textwidth}}
\toprule
页面文件 & 职责 \\
\midrule
\path{TodayScreen.tsx} & 展示当前摘要、天气、脏衣篮、频率建议和方案入口。 \\
\path{DirtyBasketScreen.tsx} & 选择或清空本次洗衣衣物，形成当前 laundry selection。 \\
\path{PlanDetailScreen.tsx} & 展示详细洗衣方案，使用对话框确认执行，并在完成后显示成功面板和撤销入口。 \\
\path{OutfitWikiScreen.tsx} & 展示穿搭分类、今日穿搭建议和加入脏衣篮操作；缺少下衣时直接提示添加裤装或裙装。 \\
\path{RecordOutfitScreen.tsx} & 记录当天穿搭，并把已穿衣物计入 wear count。 \\
\path{WardrobeScreen.tsx} & 展示衣柜列表，支持详情、删除和清空。 \\
\path{AddClothingScreen.tsx} & 手动新增衣物，或通过 ModelHub 图片/文字识别辅助填写；识图未配置时在按钮旁给出配置入口。 \\
\path{ClothingDetailScreen.tsx} & 展示单件衣物的材质、颜色、风险、洗护记忆和穿着次数。 \\
\path{LaundryRoomScreen.tsx} & 展示宿舍楼机器状态、设备编号、排队估计、天气、刷新状态和错误状态，并高亮方案推荐机器。 \\
\path{MachineDetailScreen.tsx} & 展示单台机器的状态、位置、价格和程序。 \\
\path{ReportScreen.tsx} & 渲染结构化报告、现在照做摘要、本周完成统计、风险提示和节约说明。 \\
\path{ProfileScreen.tsx} & 配置用户宿舍、楼层、取衣时间和烘干偏好；ModelHub 与演示数据放在高级设置中。 \\
\bottomrule
\end{longtable}

\subsection{前端 API/service 模块}

\begin{longtable}{p{0.32\textwidth} p{0.58\textwidth}}
\toprule
模块 & 职责 \\
\midrule
\path{mobileSummary.ts} & 组合移动端页面需要的完整摘要，包括衣柜、脏衣篮、完成洗衣记录、天气、校园上下文、频率建议、洗衣方案、烘干方案和报告；同时负责完成方案和撤销完成记录。 \\
\path{types.ts} & 定义移动端共享类型，语义上镜像 Python \path{backend/shared/models.py}，并包含 \texttt{CompletedLaundryRecord} 与 \texttt{CompletedLaundrySummary}。 \\
\path{userProfile.ts} & 读取和保存宿舍楼、楼层、最晚取衣时间、烘干偏好等用户配置。 \\
\path{modelHubConfig.ts} & 保存、清除、校验 ModelHub 配置，并限制模型名。目前允许 \texttt{gemini-3.1-pro-preview}。 \\
\path{modelHubRecognition.ts} & 在用户主动触发时发送图片或文字识别请求，要求结构化 JSON 输出。 \\
\path{clothingExtractor.ts} & 从手动输入或识别结果构造衣物信息，并转换为洗衣规划模块可使用的衣柜条目。 \\
\path{campusTowerDirectory.ts} & 提供页面可见的宿舍楼选项，隐藏 provider id。 \\
\path{campusMachineApi.ts} & 根据用户选择的宿舍楼构建校园机器上下文，归一化机器状态和排队估计。 \\
\path{weatherService.ts} & 读取清华附近 Open-Meteo 天气。 \\
\path{pricingRules.ts} & 提供洗衣、烘干、洗鞋程序的价格、时长和晾晒上下文规则。 \\
\path{frequencyAdvisor.ts} & 基于穿着次数、衣物类型和用户约束生成洗护频率建议。 \\
\path{laundryPlanner.ts} & 拆分洗衣桶、选择程序、分配机器、估算费用时间并生成提醒。 \\
\path{reportGenerator.ts} & 把 plan、drying plan、校园上下文和衣物信息转换为结构化报告。 \\
\path{outfitLogStore.ts} / \path{outfitWiki.ts} & 保存穿搭记录、统计搭配关系，并支持今日搭配建议。 \\
\path{photoThumbnail.ts} / \path{photoFileStorage.ts} & 生成并保存本地衣柜缩略图。 \\
\bottomrule
\end{longtable}

\subsection{本地状态与外部访问边界}

移动端把用户 profile、ModelHub 配置、衣柜、脏衣篮、完成洗衣记录、穿搭记录和本地图片引用保存在本机。普通衣柜、脏衣篮、计划、完成记录和报告流程不依赖独立后端服务。

完成洗衣记录保存在 \texttt{washmate.completedLaundryRecords}。每条记录包含完成时间、衣物 id、衣物名称、费用和时间估计、机器标签、方案摘要，以及完成前的穿着次数和洗涤次数快照。撤销时，移动端用该快照恢复计数，并把对应衣物放回脏衣篮。

只有三类行为会访问外部服务：

\begin{itemize}
  \item 用户主动点击图片识别、批量识别或文字识别，并且已经配置 ModelHub。
  \item 用户保存宿舍楼后，应用读取 CleverSchool 或 Haier 洗衣机状态。
  \item 应用刷新摘要时读取 Open-Meteo 天气。
\end{itemize}

缺少配置、机器状态、天气、识别结果或价格规则时，系统应该展示明确的未配置、不可用、未知或错误状态，而不是静默补一个不存在的默认值。

\subsection{后端参考实现}

\subsubsection{共享模型}

\path{backend/shared/models.py} 是 Python 后端唯一的共享模型层，存放 enum 和 dataclass。它定义了从衣物输入、衣物 profile、衣柜 item、洗衣约束、校园上下文、洗衣方案、烘干方案到报告的完整数据语言。

\begin{longtable}{p{0.26\textwidth} p{0.64\textwidth}}
\toprule
契约分组 & 主要模型 \\
\midrule
洗护和风险 & \texttt{RiskLevel}, \texttt{WashMethod}, \texttt{DryMethod}. \\
机器上下文 & \texttt{MachineType}, \texttt{MachineStatus}, \texttt{MachineTower}, \texttt{MachineInfo}, \texttt{MachineQueueEstimate}, \texttt{CampusContext}. \\
衣物和衣柜 & \texttt{ClothingInput}, \texttt{LLMResponse}, \texttt{ClothingProfile}, \texttt{WashRecord}, \texttt{WardrobeItem}. \\
规划 & \texttt{LaundryConstraints}, \texttt{FrequencyAdvice}, \texttt{LaundryChargeLine}, \texttt{LaundryBucket}, \texttt{LaundryPlan}, \texttt{DryingStep}, \texttt{DryingPlan}. \\
报告 & \texttt{WashReport}. \\
\bottomrule
\end{longtable}

\subsubsection{衣物信息抽取}

\path{backend/clothing_extraction/} 负责把用户输入、商品信息、吊牌/OCR 文本和可选图片识别结果转换为结构化 \texttt{ClothingProfile}。

\begin{itemize}
  \item \path{llm_client.py} 负责 ModelHub/Gemini v1beta 请求、图片内容、提示词和 JSON 结果格式。
  \item \path{product_info.py} 负责商品名、店铺名、吊牌文字、OCR、用户描述和补充来源归一化。
  \item \path{extractor.py} 负责解析模型输出、同步洗护动作、归一化材质比例、颜色、洗护标签、洗护符号、风险和缺失字段。
\end{itemize}

图片识别只发起一次带 JSON 格式约束的请求，使用 \texttt{responseMimeType=application/json} 和 \texttt{responseSchema}。结果包含图片类型、可见证据、材质、颜色、洗护动作、风险和缺失字段。该模块不负责衣柜存储、洗衣频率、校园机器、最终计划或报告。

\subsubsection{衣柜记忆与洗护频率}

\path{backend/wardrobe/store.py} 提供 \texttt{WardrobeStore}，负责衣柜 CRUD、洗涤历史、穿着次数、历史问题、输入校验和 JSON 序列化。\path{backend/wardrobe/frequency_advisor.py} 根据衣物、穿着次数和 \texttt{LaundryConstraints} 输出 \texttt{FrequencyAdvice}。

该模块的边界很清晰：它不调用模型，不读机器状态，也不生成最终洗衣分桶方案。

\subsubsection{校园上下文}

\path{backend/campus/machine_api.py} 提供 \texttt{LaundryMachineClient}，负责 CleverSchool 和 Haier 机器来源的请求与解析。它归一化宿舍楼、provider keys、机器 id、机器类型、状态、楼层、剩余时间和 provider 信息。

\path{backend/campus/context.py} 负责构建 \texttt{CampusContext}，把机器列表、天气、晾晒上下文和价格规则组合起来，并按机器类型生成 \texttt{MachineQueueEstimate}。如果用户选择的宿舍楼无法解析，或者调用方直接传外部楼栋 id 但没有说明 provider，该层应显式报错。

\path{backend/campus/weather.py} 负责读取并校验 Open-Meteo 天气数据。

\subsubsection{洗衣规划}

\path{backend/laundry/planner.py} 是规则化洗衣规划模块的参考实现。它负责：

\begin{itemize}
  \item 校验本次选择的衣物和 urgent item。
  \item 按不可水洗、干洗、手洗、床品、深色或掉色风险、浅色标准桶、低风险混色桶等规则拆分。
  \item 根据 \texttt{CampusContext.available\_machines} 分配洗衣机、烘干机或洗鞋机。
  \item 使用 \texttt{CampusContext.pricing\_rules} 中的真实价格和时长计算 cost breakdown。
  \item 生成洗衣液用量、洗衣袋建议、防串色、防缩水、卫生和等待时间 warning。
  \item 在用户设置预算或最大等待时间时，在 \texttt{global\_warnings} 中明确说明超预算、等待未知或等待超时。
\end{itemize}

\path{backend/laundry/load_estimator.py} 提供衣物负载估算和按负载拆分能力。\texttt{recommend\_drying} 把烘干分配从洗涤分桶中拆开，使 wash plan 和 drying plan 的责任更清晰。

\subsubsection{报告生成}

\path{backend/reports/generator.py} 把 \texttt{LaundryPlan}、可选 \texttt{DryingPlan}、衣物列表和 \texttt{CampusContext} 转换为 \texttt{WashReport}。报告模块只解释洗衣规划模块已经生成的结果，不重新分桶、不重新决策、不抓取机器、不调用模型。

\texttt{WashReport} 保留 \texttt{sections}、\texttt{action\_steps}、\texttt{cost\_breakdown}、\texttt{savings\_notes} 和 \texttt{risk\_notes} 等结构化字段，前端可以直接渲染，不需要从自然语言报告中反向解析费用和步骤。

\subsection{数据结构与主流程}

\subsubsection{数据结构对齐}

Python 和 TypeScript 分别有自己的共享契约层：

\begin{itemize}
  \item Python：\path{backend/shared/models.py}
  \item TypeScript：\path{frontend/src/api/types.ts}
\end{itemize}

TypeScript 数据结构保证 APK 内置业务服务可以独立运行。Python 数据结构保证后端参考实现和单元测试有稳定的数据语言。两者不要求文件内容逐字相同，但应保持业务语义一致。

\subsubsection{核心业务数据流}

\begin{enumerate}
  \item 用户在 Profile 保存宿舍楼、楼层、取衣时间、烘干偏好和 ModelHub 配置。
  \item 用户在 Add Clothing 手动录入衣物，或通过 ModelHub 识别辅助填写。
  \item 衣物保存为本地 wardrobe item，包含材质、颜色、风险、洗护记忆和本地图片引用。
  \item 用户在 Dirty Basket 中选择本次要洗的 item ids。
  \item \path{mobileSummary.ts} 读取衣柜、脏衣篮、天气、机器和价格规则，生成统一 mobile summary。
  \item \path{frequencyAdvisor.ts} 根据穿着次数和约束生成频率建议。
  \item \path{laundryPlanner.ts} 根据选择衣物、校园上下文和价格规则输出 \texttt{LaundryPlan}。
  \item 用户在 Plan Detail 确认执行后，\path{mobileSummary.ts} 写入完成洗衣记录，重置已洗衣物的穿着次数，增加洗涤次数，并清空本次脏衣篮。
  \item Report 页面直接渲染当前方案报告；如果本次脏衣篮已经清空，则使用完成洗衣记录展示本周统计和最近一次洗衣。
\end{enumerate}

\section{外部服务、安全与配置}

\subsection{外部服务边界}

\begin{longtable}{p{0.22\textwidth} p{0.30\textwidth} p{0.38\textwidth}}
\toprule
服务 & 触发条件 & 用途 \\
\midrule
ModelHub / Gemini v1beta & 用户配置凭据后主动点击图片识别、批量识别或文字识别。 & 把图片或文字证据转换为结构化衣物字段。 \\
CleverSchool & 用户保存支持的宿舍楼。 & 读取智慧校园洗衣机楼栋和机器状态。 \\
Haier & 用户保存支持的宿舍楼。 & 读取海乐生活点位和设备状态。 \\
Open-Meteo & 刷新 mobile summary。 & 读取清华附近天气，用于晾晒上下文和页面展示。 \\
\bottomrule
\end{longtable}

\subsection{在线数据来源与字段}

天气数据由移动端 \path{frontend/src/api/weatherService.ts} 和 Python 参考实现 \path{backend/campus/weather.py} 读取。当前接口是 Open-Meteo forecast API：\url{https://api.open-meteo.com/v1/forecast}。移动端使用清华附近坐标 \texttt{40.0023, 116.3268}，时区为 \texttt{Asia/Shanghai}，请求当前天气字段 \path{temperature_2m}、\path{relative_humidity_2m}、\path{precipitation} 和 \path{weather_code}。返回结果必须包含 \path{current}。温度、湿度和降水必须是有效数字。读取成功时，页面收到 \texttt{source=open-meteo} 和 \texttt{status=live}；读取失败或字段不合法时，页面收到 \texttt{status=unavailable} 和错误信息。

洗衣机数据来自两类校园服务。移动端先使用 \path{frontend/src/api/campusTowerDirectory.ts} 中的本地宿舍楼目录，把用户看到的宿舍楼名映射到内部 \path{towerKey} 或 \path{positionId}。页面不会直接展示这些内部 id。

CleverSchool 机器状态由移动端 \path{frontend/src/api/campusMachineApi.ts} 读取。接口地址是 \url{https://api.cleverschool.cn/washapi4/device/status}。请求使用 POST，body 包含 \path{towerKey} 和空的 \path{deviceType}。响应中的 \path{macUnionCode} 用来拆出机器编号，\path{tower} 和 \path{floorName} 用来生成位置，\path{status} 用来生成机器状态和剩余时间。状态文字中的“待机”“空闲”“可用”“可使用”会转成 \texttt{available}；“工作”“运转”“使用”“运行”会转成 \texttt{running}；“脱水”“开盖”“出错”“错误”“异常”“故障”“不可使用”“停用”会转成 \texttt{out\_of\_service}。剩余时间只从明确的“剩余时间”文字中读取。

海乐生活机器状态也由 \path{campusMachineApi.ts} 读取。接口地址是 \url{https://yshz-user.haier-ioc.com/position/deviceDetailPage}。请求使用 POST，body 包含 \path{positionId}、\path{categoryCode}、\path{page}、\path{floorCode} 和 \path{pageSize}。移动端会分别请求 \texttt{00}、\texttt{01} 和 \texttt{02} 三类设备，并把它们映射为普通洗衣机、洗鞋机和烘干机。响应中的 \path{id} 是机器编号，\path{name} 是位置文字，\path{state} 是状态。状态 \texttt{1} 表示可用，\texttt{2} 表示运行中，\texttt{3} 表示不可用；其他值记为 \texttt{unknown}。

Python 参考实现还保留楼栋和点位发现接口。CleverSchool 使用 \url{https://api.cleverschool.cn/washapi4/device/tower} 获取楼栋列表。海乐生活使用 \url{https://yshz-user.haier-ioc.com/position/nearPosition}，并用清华附近坐标 \texttt{lng=116.32697, lat=40.00281} 获取附近点位。后端会把同名楼栋合并成 \texttt{MachineTower}，并在 \path{provider_keys} 中保留不同来源的内部 id。

实时机器接口不提供完整价格和程序时，系统不会在机器解析层补齐这些业务数据。价格、模式和时长来自 \path{config/machine_rules.json} 以及移动端镜像的 \path{frontend/src/api/pricingRules.ts}。本地网页预览时，请求会改写到 \path{/cleverschool-api} 和 \path{/haier-api} 代理路径；Android APK 中使用 Capacitor HTTP 直接发起请求。

\subsection{密钥与本地配置}

移动端 release APK 不内置真实 API key。用户在 Profile 中填写 ModelHub \texttt{baseUrl}、\texttt{apikey} 和 \texttt{model\_name}，配置保存在本机，并提供清除入口。

Python 参考实现使用本地 \path{config/api_config.json}，仓库只提交 \path{config/api_config.example.json}。真实 API key、\path{.env}、keystore、签名密码和真实用户数据都不应提交到仓库。

\subsection{不静默猜测业务数据}

项目规则强调 no fallback / no default。这里的含义是：缺少业务数据时，应显式展示未配置、不可用、未知、错误状态或抛出异常，而不是补一个看似合理但没有来源的数据。不能凭空生成材质、颜色、机器状态、天气、价格、时长、风险或洗衣方案。

\section{运行方法与交付文件}

\subsection{本地运行与构建命令}

项目仓库地址为 \url{https://github.com/vegetable123q/Wash-agent}。debug APK 交付路径为 \path{frontend/android/app/build/outputs/apk/debug/app-debug.apk}。如果本地尚未生成 APK，可按下面命令重新构建。

\begin{lstlisting}[style=command]
uv sync
uv run python -m unittest discover -v

cd frontend
npm install
npm test
npm run build
npm run cap:sync
\end{lstlisting}

构建 debug APK：

\begin{lstlisting}[style=command]
cd frontend
npm run apk:debug
\end{lstlisting}

生成路径：

\begin{lstlisting}[style=command]
frontend/android/app/build/outputs/apk/debug/app-debug.apk
\end{lstlisting}

\subsection{自动检查与发布}

\path{.github/workflows/preview.yml} 用于 preview 检查。它会同步 Python 环境，运行 Python 单元测试，使用 \texttt{npm ci} 安装前端依赖，运行前端测试，并执行前端 build。

\path{.github/workflows/release-apk.yml} 用于 main 分支发布。它会运行 Python 和前端测试，从 \path{frontend/package.json} 读取 release version，更新 Android version，安装 JDK 21 和 Android SDK 35，执行 Capacitor sync，构建签名 release APK，上传 artifact，并发布 GitHub Release。

稳定 release 签名依赖 GitHub Secrets：\texttt{ANDROID\_KEYSTORE\_BASE64}、\texttt{ANDROID\_KEYSTORE\_PASSWORD}、\texttt{ANDROID\_KEY\_ALIAS}、\texttt{ANDROID\_KEY\_PASSWORD}。如果这些 secrets 缺失，CI 可以生成临时 key 来构建可安装 APK，但该 key 不适合作为长期更新签名。

\section{测试与验收结果}

\subsection{自动化测试与覆盖}

当前仓库包含 Python unittest 和 frontend Vitest。后端测试覆盖衣物抽取、衣柜 CRUD、频率建议、校园机器解析、校园上下文、天气、洗衣规划、报告生成和完整集成流程。前端测试覆盖 API services、页面渲染、Profile 配置、ModelHub 配置持久化、脏衣篮、完成洗衣记录与撤销、推荐机器高亮、报告周统计、刷新行为和 App integration。

\subsection{实测验收结果}

\path{evaluation_system.md} 记录了 2026-05-28 的验收结果。主要结论如下：

\begin{itemize}
  \item Python unittest：通过，命令为 \texttt{uv run python -m unittest discover -v}。
  \item Frontend test：通过，命令为 \texttt{npm test -- --run}。
  \item Frontend build：通过，命令为 \texttt{npm run build}。
  \item Capacitor sync：通过，命令为 \texttt{npm run cap:sync}。
  \item APK debug build：通过，命令为 \texttt{npm run apk:debug}，生成 \path{frontend/android/app/build/outputs/apk/debug/app-debug.apk}。
  \item APK 安装与启动：通过，Android Emulator 安装后可打开 \texttt{edu.tsinghua.washmatecampus/.MainActivity}。
  \item APK 内置服务：通过，不启动本机后端也可以查看 Today、衣柜、洗衣房、方案和报告等主流程。
  \item WebView 冒烟测试：通过，覆盖 Today、衣柜、衣物详情、添加衣物、洗衣房、机器详情、方案详情、报告和我的页面。
\end{itemize}

本轮验收发现的 P1 问题已修复并复验，包括衣柜删除、宿舍楼选择、图片识别配置、展示型按钮禁用、空衣柜状态、固定样例详情、APK 后端依赖、API key 配置和等待时间提醒等。当前交付结论为：自动化、模拟器 APK、API 和移动端主流程均通过，未发现剩余 P0/P1 缺陷。

\section{团队分工}

\subsection{成员与贡献预留}

以下表格用于填写最终团队分工和贡献说明。成员名单已按提交信息预置，具体负责模块和完成内容可在最终版中补充。

\begin{longtable}{p{0.18\textwidth} p{0.30\textwidth} p{0.42\textwidth}}
\toprule
成员 & 负责模块 & 主要贡献 \\
\midrule
武子皓 & 待补充 & 待补充 \\
王成思宇 & 待补充 & 待补充 \\
林恺鋆 & 待补充 & 待补充 \\
王博翔 & 待补充 & 待补充 \\
周逸航 & 待补充 & 待补充 \\
\bottomrule
\end{longtable}

\subsection{模块维护说明}

\begin{itemize}
  \item 页面和交互放在 \path{frontend/src/screens/}，通用外壳或组件放在 \path{frontend/src/components/}。
  \item 移动端业务逻辑放在对应 \path{frontend/src/api/} service 中。
  \item Python 参考逻辑放在对应 backend package 中，例如 clothing extraction、wardrobe、campus、laundry 或 reports。
  \item Python 共享模型只放在 \path{backend/shared/models.py}。
  \item 移动端共享类型放在 \path{frontend/src/api/types.ts}。
  \item 如果修改模块边界、公共接口、配置行为或主流程，应同步更新 \path{docs/structure_design.md}。
\end{itemize}

\subsection{开发注意事项}

\begin{itemize}
  \item 不新增旧路径兼容层，例如重新创建 flat backend import。
  \item 不把提示词、解析器、洗衣规则、存储细节或报告拼装写进页面。
  \item 不为普通衣柜、脏衣篮、方案和报告流程引入隐藏的远程后端依赖。
  \item 不提交 API key、keystore、签名密码或真实用户数据。
  \item 不通过猜测结果掩盖外部 API 失败或缺失配置。
\end{itemize}

\section{当前限制与未来改进}

当前实现已经适合作为完整移动端产品演示，但仍有明确扩展空间：

\begin{itemize}
  \item 账号体系和云端同步尚未实现；当前衣柜、profile、脏衣篮和完成洗衣记录以本机存储为主。
  \item 如果未来要支持多设备同步，可以增加 hosted backend，但不应破坏当前 APK 内置业务服务的清晰边界。
  \item 如果要支持更多 ModelHub 模型，应通过显式 allow-list、UI 控件和测试来扩展。
  \item 如果要支持更多校园机器 provider，应放在 campus API/context 边界内，并继续向 UI 暴露用户可读楼名而不是 provider id。
  \item 如果要导出更多报告格式，应消费结构化 \texttt{WashReport}，不要解析自然语言段落。
\end{itemize}

\subsection{总结}

WashMate Campus 的事实实现是一个以移动端为主的完整产品。前端不仅是 UI，也包含 APK 内置业务服务；Python 后端是参考实现和测试基准，帮助团队保持清晰的模块边界、数据结构和业务规则。这样的设计让项目能直接演示完整产品流程，也能通过测试和文档保留可维护的工程结构。

\end{document}
